% ステップ0 検証用: 色付き版（ラベルの正解に相当）
% probe_plain.tex と本文はまったく同じ。色の指定だけを足してある。
\documentclass[uplatex,dvipdfmx]{jsarticle}
\usepackage{xcolor}

% ラベルごとの色。ai/auto_label.py の COLOR_TO_LABEL と値を一致させること。
\definecolor{labelheading}{HTML}{FF0000}  % heading
\definecolor{labelcaption}{HTML}{00AA00}  % caption
\definecolor{labelother}{HTML}{888888}    % other

\renewcommand{\thepage}{\textcolor{labelother}{\arabic{page}}}
\begin{document}

\section{\color{labelheading}はじめに}

本研究では、ブロックエディタで編集した内容をコンパイル可能なTeXコードとして
出力するツールを開発する。本節ではその背景と目的について述べる。
既存のツールでは、TeXの文法を習得していない利用者が文書を作成することが難しい。

\subsection{\color{labelheading}目的}

PDFを読み込み、自作Transformerによるブロック分類を行い、
ブロックエディタ上で編集可能な形に変換することを目的とする。

\begin{figure}[h]
\centering
\fbox{\rule{0pt}{3cm}\rule{5cm}{0pt}}
\caption{\color{labelcaption}システム全体の構成}
\end{figure}

\begin{table}[h]
\centering
\caption{\color{labelcaption}対応ブロック種別}
\begin{tabular}{|l|l|}
\hline
種別 & 対応フェーズ \\
\hline
見出し & Phase 1 \\
段落 & Phase 1 \\
表 & Phase 2 \\
\hline
\end{tabular}
\end{table}

\end{document}